%%%%%%%%%%%%%%%%%%%%%%%%%%%%%%%%%%%%%%%%%%%%%%%%%%%%%%%%%%%%%%%%%%%%%%
% Lớp 8 | Chương X. Một số hình khối trong thực tiễn | Hình chóp tam giác đều
% Dạng 8H0?38-3: Bài toán thực tế
%%%%%%%%%%%%%%%%%%%%%%%%%%%%%%%%%%%%%%%%%%%%%%%%%%%%%%%%%%%%%%%%%%%%%%


% nguồn: De_AI_Toan8_cuoi_hoc_ki_i_1 (nhập 21/11/2026)
\begin{bt}%[8H0N38-3]%[ma:50ba8a33]
	Một khối đồ chơi Rubik có dạng hình chóp tam giác đều với diện tích mặt đáy khoảng $28\text{ cm}^2$ và chiều cao bằng $9\text{ cm}$. Tính thể tích của khối Rubik đó.
	\loigiai{
		Thể tích của khối Rubik dạng hình chóp tam giác đều là: \\
$V = \dfrac{1}{3} S_{\text{đáy}} \cdot h = \dfrac{1}{3} \cdot 28 \cdot 9 = 84\text{ cm}^3$. \\
Vậy thể tích của khối Rubik đó là $84\text{ cm}^3$.
	}
\end{bt}

% nguồn: De_AI_Toan8_cuoi_hoc_ki_i_1 (nhập 21/11/2026)
\begin{bt}%[8H0N38-3]%[ma:70c144d4]
	Bạn Lan làm một chiếc đèn lồng có dạng hình chóp tam giác đều với độ dài cạnh đáy $16\text{ cm}$ và độ dài trung đoạn $15\text{ cm}$. Tính diện tích giấy màu tối thiểu cần dùng để dán kín ba mặt bên của chiếc đèn lồng (coi phần mép dán không đáng kể).
	\loigiai{
		Diện tích giấy màu cần dùng chính là diện tích xung quanh của hình chóp tam giác đều. \\
Nửa chu vi đáy là $p = \dfrac{3 \cdot 16}{2} = 24\text{ cm}$. \\
Diện tích xung quanh là $S_{xq} = p \cdot d = 24 \cdot 15 = 360\text{ cm}^2$. \\
Vậy bạn Lan cần ít nhất $360\text{ cm}^2$ giấy màu.
	}
\end{bt}

% nguồn: De_AI_Toan8_cuoi_hoc_ki_i_1 (nhập 21/11/2026)
\begin{bt}%[8H0N38-3]%[ma:f1bda04e]
	Một khối pha lê lưu niệm có dạng hình chóp tứ giác đều với cạnh đáy bằng $9\text{ cm}$ và chiều cao bằng $10\text{ cm}$. Tính thể tích của khối pha lê đó.
	\loigiai{
		Diện tích mặt đáy hình vuông của khối pha lê là $S_{\text{đáy}} = 9^2 = 81\text{ cm}^2$. \\
Thể tích của khối pha lê là: \\
$V = \dfrac{1}{3} S_{\text{đáy}} \cdot h = \dfrac{1}{3} \cdot 81 \cdot 10 = 270\text{ cm}^3$. \\
Vậy thể tích của khối pha lê là $270\text{ cm}^3$.
	}
\end{bt}
